\chapter{Cluster Transition Screening}\label{ch:cluster_transition_screening}

\section{Introduction}\label{sec:ch5_introduction}

Genomic surveillance can reveal many groups of genetically similar infections, but deciding which groups warrant investigation requires more than counting their members. Chapter~\ref{ch:genomic_networks} showed that compatibility was geographically structured and that graph size and connectivity varied with sequencing volume and case counts. A large cluster may therefore reflect a period of extensive sampling, and concentration within one residential area may be common in its surveillance context. This chapter asks which clusters stand out relative to other observed clusters, and whether those with large accumulations of sampled infections are also associated with extensive later sampled representation.

There are epidemiological reasons to consider these patterns separately. In Boston, an outbreak in a skilled nursing facility involved many infections but little wider spread, whereas a conference outbreak was followed by extensive community and international spread \citep{Lemieux2021}. This contrast motivates examining both the concentration of sampled infections and their later representation. It does not establish that either pattern, when observed without exposure information, identifies a superspreading event. As Chapter~\ref{ch:literature} explains, the Scottish clusters do not supply the individual secondary-infection counts needed for that classification.

We use three observable features. Size relative to other clusters helps distinguish a large accumulation from the background in which it was sampled. Membership absent from earlier linked clusters distinguishes newly represented sequences from membership already present in the preceding window. Additional sequences in later linked clusters describe subsequent sampled representation, considered relative to the starting cluster's size. These features motivate two screening summaries: the \emph{local burst score} combines relative size and available membership novelty, while the \emph{onward burden score} combines direct and cumulative later representation. Here, `local' refers to a cluster and its immediate links, rather than a geographic area, and `burden' counts sampled sequences rather than illness severity.

To obtain these features, we connect clusters in successive surveillance windows when they share sampled sequences. In the compatibility graph, each node represents a sequence; in this \emph{transition graph}, each node represents a cluster in a given window. We refer to the latter as clusters throughout this chapter. The links describe shared membership, without establishing transmission between clusters. We first describe the clusters, then compare their two scores with permutation reference distributions and examine sensitivity to screening thresholds. Eligible clusters meeting either criterion receive candidate superspreading-compatible (CSC) status for follow-up investigation.

\section{Materials and Methods}\label{sec:ch5_methods}

We analysed the high-quality Scottish SARS-CoV-2 genomic clustering dataset constructed in Chapter~\ref{ch:genomic_networks}. Each record represents one sequence in one 21-day surveillance window, with its cluster assignment and information about the epidemic and sampling context. Clusters were defined separately within each window and lineage. The same sequence can contribute records to overlapping windows, and clusters linked across windows need not retain the same membership.

We used the relatively coarse communities obtained at the primary Leiden resolution ($\gamma = 0.3$). Descriptive summaries used all 134 rolling 21-day surveillance windows, which advance in 7-day steps and overlap by 14 days. The transition graph used 67 windows, retaining every other source window (Appendix~Figure~\ref{fig:alternate_window_retention}). The resulting 14-day step and seven-day overlap allow a sequence to appear in at most two consecutive retained windows, rather than up to three source windows. This choice reduces repeated inclusion while retaining an interval in which shared sequences can connect successive cluster assignments. A design with no overlap would prevent links based on shared sequences. The selected overlap is an operational trade-off; these analyses do not establish an optimal interval or equate the window step with a fixed number of transmission generations.

The remainder of this section describes how we evaluate genetic and temporal cohesion within clusters (Section~\ref{sec:ch5_pairwise_distance_summary}), construct a directed transition graph linking clusters in adjacent retained windows (Section~\ref{sec:ch5_transition_graph}), and apply a two-dimensional screening procedure to assign candidate labels (Section~\ref{sec:ch5_screening_score}). Appendix~\ref{app:cluster_transition_screening} provides further details on graph construction, screening scores, and permutation calibration.

\FloatBarrier

\begin{tcolorbox}[
    title=Glossary: cluster summaries and links,
    halign=left,
    fontupper=\small,
    ]
    \begin{description}
        \item[Cluster:] A group of sequences assigned to the same Leiden community within one window and lineage. Counts across windows count these separate cluster assignments, not distinct outbreaks.
        \item[Size and duration:] Size counts unique sequences in a cluster. Duration is the number of days between its earliest and latest collection dates.
        \item[Residential spread and distance:] Spread counts represented Data Zones. Distance is the median pairwise distance between their residential centroids, in kilometres. These describe different aspects of geography.
        \item[Transition link:] A connection between clusters in consecutive retained windows that share at least one sequence; its weight counts the shared sequences.
        \item[Weak component:] A set of clusters connected through transition links when link directions are ignored, with no links to other clusters outside the set.
        \item[Membership novelty:] The fraction of a cluster's sequences absent from its directly linked earlier clusters; termed upstream novelty in the score formula.
        \item[Downstream burden:] Additional sampled sequences in directly linked later clusters, or in all clusters reachable by following links forward. It describes sampled membership rather than secondary infections.
    \end{description}
\end{tcolorbox}

\subsection{Cluster genetic and temporal cohesion}\label{sec:ch5_pairwise_distance_summary}

We described genetic and temporal distances within the constructed clusters. Genetic differences and collection dates already contribute to compatibility scores, so these summaries describe the resulting membership rather than independently validate transmission links. We extracted SNP (single-nucleotide polymorphism) and collection-date differences for all pairs assigned to the same cluster, grouping the pairs by surveillance window and lineage. Groups containing fewer than ten within-cluster pairs were excluded. We summarised the group medians with equal weighting and with weights proportional to the number of pairs. Appendix~Table~\ref{tab:cluster_pairwise_distance_summary} provides the full distributions and weighting details.

\subsection{Construction of cluster transition graph}\label{sec:ch5_transition_graph}

We constructed a directed transition graph by linking clusters in adjacent retained windows that share sampled sequences. Because retained windows overlap, the same sequence may appear in two consecutive retained windows. This shared membership links cluster assignments across consecutive retained windows, but it does not imply direct transmission between individuals or outbreaks.

\begin{tcolorbox}[title=Illustrative example: three successive windows,fontupper=\small]
Consider three clusters, each containing six sequences. Letters identify distinct sampled sequences; this is an illustration, not an observed outbreak.
\begin{center}
\begin{tabular}{@{}ll@{}}
Earlier cluster & $a,b,c,d,e,f$ \\
Focal cluster & $e,f,g,h,i,j$ \\
Later cluster & $i,j,k,l,m,n$ \\
\end{tabular}
\end{center}
The focal cluster shares $e,f$ with the earlier cluster. Its other four sequences are absent from that earlier membership, giving upstream novelty of $4/6$. This denotes newly represented membership, not four newly acquired infections.

The later cluster shares $i,j$ with the focal cluster and contains four additional sequences, $k,l,m,n$. If these are the only links, direct downstream burden for the focal cluster is four. Cumulative burden is also four here; in a longer chain it includes unique additional sequences in all later reachable clusters, counting each sequence only once. These counts do not assign a screening score: percentile ranks and calibrated values require the other eligible clusters used for comparison.
\end{tcolorbox}

In the notation below, $V$ is the set of clusters and $E$ the set of links in the directed transition graph $G=(V,E)$. For cluster $v \in V$, let $S_v$ contain its unique sequences and let $t_v$ be its retained-window index. A link $u \rightarrow v$ is present when $v$ is in the next retained window and the two clusters share at least one sequence:

\nopagebreak[4]
\[
	t_v = t_u + 1 \quad \text{and} \quad S_u \cap S_v \neq \varnothing.
\]

The edge weight is

\nopagebreak[4]
\[
	w_{uv} = |S_u \cap S_v|,
\]

representing the number of unique sequences shared by the two cluster assignments. Edges advance exactly one retained window forward in time. The graph is therefore directed and acyclic: edges point forward in time and cannot form a closed loop.

We classified each cluster by the number of directly linked earlier and later clusters, termed its in-degree and out-degree (Figure~\ref{fig:transition_graph_roles}). For example, one earlier and one later link form a continuation; several later links form a branch; several earlier links form a merge. The nine combinations distinguish these patterns from clusters with no observed earlier or later link. Clusters without earlier links are termed sources, those without later links sinks, and those without either isolated. Sources in the first retained window were flagged as left-censored and sinks in the last as right-censored, because the corresponding earlier or later observations were outside the study period.

\begin{figure}[htbp]
	\centering
	\includegraphics[width=\textwidth]{fig_transition_graph_roles}
	\caption[Graph-role definitions in the transition graph]{
    \textbf{Graph-role definitions in the transition graph.} Each circle represents a cluster. A coloured focal cluster is connected to grey earlier or later clusters by shared-membership links. Rows and columns distinguish the numbers of earlier and later links, respectively, defining isolated clusters, boundaries, continuations, branches, and merges.
	}\label{fig:transition_graph_roles}
\end{figure}

The screening quantities compare a cluster's membership with earlier and later linked membership. For the formulas, directly linked earlier clusters are called parents, directly linked later clusters children, and all clusters reachable by following links forward descendants. These terms describe graph connections. For a focal cluster $v$ with membership $S_v$, let $\mathcal{P}_v=\{p:p\rightarrow v\}$ contain its directly linked earlier clusters.

When earlier links are observed, we first combine all sequences in the directly linked earlier clusters into one comparison set, counting each sequence once:

\nopagebreak[4]
\[
    S_{\mathcal{P}_v}=\bigcup_{p\in\mathcal{P}_v} S_p,
\]

Upstream novelty is the fraction of the focal cluster's sequences absent from that comparison set:

\nopagebreak[4]
\[
    U_v=\frac{|S_v\setminus S_{\mathcal{P}_v}|}{|S_v|}.
\]

For clusters without an observed earlier link, $U_v$ is unavailable. An absent link does not show that all membership is newly represented: earlier related infections may not have been sampled or linked.

We computed two downstream quantities to describe additional sampled sequences found in later linked clusters. Full formulae are provided in Appendix~\ref{app:cluster_transition_screening}.

\begin{description}
    \item[Direct downstream burden:] A weighted count of additional sequences in directly linked later clusters. If a later cluster has several earlier links, its contribution to the focal cluster is weighted by the focal link's share of its total incoming shared membership. The formula in Appendix~\ref{app:cluster_transition_screening} specifies this allocation.

    \item[Cumulative downstream burden:] The number of unique sequences absent from the focal cluster but present in any cluster reachable by following links forward. A sequence is counted once even if it appears in several later clusters.
\end{description}

These quantities describe additional sampled sequences in the transition graph. They do not establish which infections gave rise to later infections, count secondary infections, or measure transmission from one cluster to another. Clusters with no positive direct or cumulative downstream burden were excluded from onward burden scoring. This means that no additional sampled sequences were observed in later linked clusters; it does not establish that transmission ended.

Upstream novelty and the two downstream burden quantities are used to construct the local burst and onward burden screening scores described below.

\subsection{Screening score construction and candidate classification}\label{sec:ch5_screening_score}

We screened clusters meeting a minimum size threshold. This excluded the smallest groups, in which a difference of one sequence can substantially change the membership fractions. Sensitivity analyses evaluated thresholds from 2 to 20 sequences in increments of two, recomputing scores and permutation calibration for each eligible set. These analyses describe how the operational choice changes selection; they do not establish a size threshold that identifies true superspreading events.

We first compared cluster size with the most specific suitable group: screened clusters from the same window and clade, then the same window, then all screened clusters. A group required at least 20 clusters and non-zero variance. The \emph{contextual excess size}, $Z_v$, expresses log-transformed size relative to that group's mean and standard deviation (Appendix~\ref{app:cluster_transition_screening}). This locates a cluster within an observed comparison group; it does not estimate how large the cluster would have been under complete sampling.

We converted components to percentile ranks within each retained window. A high rank means that the component is large relative to the other applicable clusters in that window. Ranking places components on a common scale and limits the influence of extreme magnitudes when they are combined. Equal weighting then gives the two available ranks the same contribution. Ranking, equal weighting, and the size adjustment below are operational design choices, not weights estimated from confirmed outbreaks.

\subsubsection*{Local burst score}

The local burst score (LBS)\nomenclature{LBS}{local burst score} combines contextual excess size with upstream novelty. Let $R_{v,Z}$ be the within-window percentile rank of $Z_v$ among screened clusters, and let $R_{v,U}$ be the corresponding rank of $U_v$ among clusters for which upstream novelty is observable. We define

\nopagebreak[4]
\[
    LBS_v =
    \begin{dcases}
        \frac{1}{2}(R_{v,Z} + R_{v,U}) & \text{if } U_v \text{ is observable}, \\
        R_{v,Z} & \text{otherwise}.
    \end{dcases}
\]

Clusters with earlier links are scored by both relative size and membership novelty. Clusters without earlier links are scored by size alone because their novelty cannot be observed. Since the score averages ranks, a high value need not mean that both components exceed a separate threshold.

\subsubsection*{Onward burden score}

The onward burden score (OBS)\nomenclature{OBS}{onward burden score} combines two measures expressed relative to the starting cluster's size: direct downstream burden and cumulative downstream burden. For each cluster with positive observed direct or cumulative downstream burden (a burden-eligible cluster), we adjusted both quantities for cluster size and then ranked them within the retained window (Appendix~\ref{app:cluster_transition_screening}). Let $R_{v,A}$ and $R_{v,C}$ be the corresponding within-window percentile ranks among burden-eligible clusters. The onward burden score is

\nopagebreak[4]
\[
    OBS_v = \frac{1}{2} \left( R_{v,A} + R_{v,C} \right).
\]

OBS is defined only for screened clusters with positive observed direct or cumulative downstream burden. Clusters with no positive observed downstream burden are marked as not applicable for burden scoring rather than given ranks of zero.

The size adjustment asks whether later representation is large relative to the starting membership: the same additional count contributes a larger adjusted value for a smaller starting cluster. Direct burden describes the next retained window, while cumulative burden includes more distant linked membership. Equal weighting retains both perspectives, although the two quantities overlap and are not independent measures of transmission.

\subsubsection*{Permutation calibration}

Calibration compares an observed score with scores obtained by reassigning observed component profiles among clusters in a specified comparison group, termed a calibration stratum. A profile contains a cluster's component values and any unavailable values. In each of 1,000 permutations, profiles were moved together rather than shuffling their components separately; percentile ranks within windows and composite scores were then recomputed. The reference therefore retains observed combinations of size, novelty, and burden. It does not generate an epidemic without superspreading or simulate unsampled infections.

For local burst, comparison groups first separated clusters according to whether upstream novelty was available, then used window and clade where at least 20 clusters were available. Broader groups were used otherwise. For onward burden, permutation was restricted to burden-eligible clusters and grouped by clade where available. Burden profiles could therefore move between windows within a clade, while their percentile ranks were recomputed within the receiving window.

For each score, an upper-tail $p$-value compares the observed value with its permutation reference (Appendix~\ref{sec:app_ch5_permutation_calibration}). Small values identify high scores relative to this reference. For cluster $v$ and score label $k\in\{\mathrm{LBS},\mathrm{OBS}\}$, the conservative value $p_{v,k}^{\mathrm{cons}}$ counts every reference score that exceeds or ties the observed score; the randomised value $p_{v,k}^{\mathrm{rand}}$ counts every exceedance and a randomly chosen fraction of exact ties. The add-one correction prevents zero values. We used randomised values for classification and retained conservative values for audit. Interpreting these as inferential $p$-values requires the assumed reassignment of profiles to be valid under the relevant null hypothesis \citep{Hemerik2017}; the reference construction alone does not establish that assumption.

\subsubsection*{Candidate classification}

Candidate classification used the primary specification selected after the threshold sensitivity analysis reported in the Results: $n_v \geq 6$ and $p_{v,k}^{\mathrm{rand}} \leq 0.05$ for at least one available score $k$. The sensitivity grid also considered calibrated upper-tail thresholds ($p_{v,k}^{\mathrm{rand}}$) from 0.01 to 0.10 in increments of 0.01.

The primary cutoffs define an operational review set. The threshold applies separately to each score, and CSC status is assigned when either qualifies. It is not an estimated probability of superspreading or a demonstrated 5\% error rate for the combined candidate set.

Under the primary specification, we assigned each cluster to one mutually exclusive tier using the rules in Table~\ref{tab:candidate_tiers}. Here, $p_{v,\mathrm{LBS}}^{\mathrm{rand}}$ and $p_{v,\mathrm{OBS}}^{\mathrm{rand}}$ denote the calibrated $p$-values for the local burst and onward burden scores, when available. We defined candidate superspreading-compatible (CSC) clusters as those assigned to any of the three high-priority tiers.

\begin{table}[htbp]
    \centering
    \caption[Screening-tier definitions for screened clusters]{
        Screening-tier definitions at the primary specification. Rules below the size-ineligible row apply to
        clusters with $n_v \geq 6$, in the order shown. CSC clusters comprise the three high-priority tiers.
    }\label{tab:candidate_tiers}
    \begin{thesistablebody}[\thesistablesetup\renewcommand{\arraystretch}{1.5}]{@{}P{0.32\textwidth}P{0.63\textwidth}@{}}
        \toprule
        \textbf{Tier} & \textbf{Assignment rule} \\
        \midrule
        Size ineligible
        & $n_v < 6$ \\

        High priority, both dimensions
        & $p_{v,\mathrm{LBS}}^{\mathrm{rand}} \leq 0.05$ and $p_{v,\mathrm{OBS}}^{\mathrm{rand}} \leq 0.05$ \\

        High priority, burst only
        & $p_{v,\mathrm{LBS}}^{\mathrm{rand}} \leq 0.05$ and the burden criterion was not met \\

        High priority, burden only
        & $p_{v,\mathrm{OBS}}^{\mathrm{rand}} \leq 0.05$ and the burst criterion was not met \\

        Possible review
        & No high-priority criterion was met, but at least one applicable score had $p_{v,k}^{\mathrm{rand}} \leq 0.10$ \\

        High score, uncalibrated
        & No high-priority criterion was met, and at least one raw score was $\geq 0.90$ with unavailable calibration \\

        Background or low information
        & None of the criteria above were met \\
        \bottomrule
    \end{thesistablebody}
\end{table}

Together, this procedure yields a directed transition graph in which each size-eligible cluster has calibrated local burst and onward burden measures, where applicable, and a single assigned screening tier. The resulting CSC clusters at the primary specification are characterised in Chapter~\ref{ch:candidate_characterisation}.

\section{Results}\label{sec:ch5_results}

The pipeline in Chapter~\ref{ch:genomic_networks} identified approximately 551,000 clusters across 134 surveillance windows and 790 Pango lineages. These are separate cluster assignments within windows, not counts of distinct outbreaks. Most clusters (82.5\%) contained one sequence. Within-cluster distance summaries were available for 1,707 eligible combinations of window and lineage, with medians of 0 SNPs and 4 days under both weighting methods (Table~\ref{tab:cluster_pairwise_distance_summary}). These values describe the membership produced by the compatibility and clustering procedure.

We present the remaining results in three parts. First, we describe cluster size, duration, residential spread, and residential distance (Glossary above). Second, we show how clusters connect through time in the transition graph. Third, we report the screening outcomes.

\subsection{Cluster size, timing, and geographic spread}\label{sec:ch5_cluster_landscape}

Most clusters were small, while a few were much larger (Figure~\ref{fig:cluster_landscape}A). Across policy periods, the median non-singleton cluster contained 2--3 sequences, but the largest cluster sizes varied substantially (Table~\ref{tab:cluster_period_summary}). The largest clusters occurred during high-incidence Omicron periods, and the maximum cluster size differed more than tenfold between windows despite stable median size (Figure~\ref{fig:cluster_landscape}B).

\begin{figure}[htbp]
	\centering
	\includegraphics[width=\textwidth]{fig_cluster_landscape}
	\caption[Scottish EpiLink cluster landscape]{
		\textbf{Window-lineage EpiLink cluster landscape.} (A) Complementary cumulative distribution of cluster size across all window-lineage clusters (log--log axes). (B) Median, 90th percentile, and maximum non-singleton cluster size across rolling windows (log $y$-axis); shaded bands denote epidemic eras. (C) Maximum cluster size in each window against window sequencing coverage. (D) Joint distribution of non-singleton cluster size and median pairwise residential Data Zone-centroid distance (log--log hexbin); colour indicates the number of clusters in each bin. (E) Window-level Pearson correlations of positive-test counts, sequence counts, and sequencing coverage with cluster count and size summaries. (F) Window-level Pearson correlations between cluster-size summaries and temporal, Data Zone, and residential-distance summaries. %chktex 8
		}\label{fig:cluster_landscape}
\end{figure}

\input{assets/scotland/tab_cluster_period_summary}

Positive-test counts, sequence counts, and sequencing coverage were associated with different parts of the cluster-size distribution (Figure~\ref{fig:cluster_landscape}E). Positive-test counts and sequence volumes were strongly associated with cluster counts and maximum cluster size, but only weakly associated with typical cluster size. Sequencing coverage showed the opposite pattern: it was weakly associated with cluster counts but more clearly related to the upper part of the non-singleton size distribution. Large cluster maxima occurred across the full range of coverage values (Figure~\ref{fig:cluster_landscape}C).

Clusters spanned short periods and few residential areas. Median duration was well below the 21-day window length, and the typical non-singleton cluster spanned only 2 Data Zones (Table~\ref{tab:cluster_period_summary}). Cluster size was closely related to the number of Data Zones represented, but not to residential distance (Figure~\ref{fig:cluster_landscape}D,F). Thus, larger clusters usually included more residential areas, but they were not necessarily more geographically dispersed.

\subsection{Topology and connectivity of the cluster transition graph}\label{sec:ch5_transition_structure}

We next linked clusters across the 67 retained windows into a directed transition graph. This graph captures continuity in sampled sequence membership and provides the structure used for local burst and onward burden screening.

The transition graph represented 277,154 clusters across the retained windows and 117,583 directed links (Table~\ref{tab:transition_graph_summary}). These counts describe cluster assignments and their shared membership across windows. Clusters had few links overall: the mean numbers of incoming and outgoing links were both 0.42. Cluster and link counts varied strongly across time, peaking during high-volume Omicron windows (Figure~\ref{fig:transition_graph_characteristics}A).

\begin{table}[htbp]
	\centering
	\caption[Summary of transition-graph structure]{
		Summary of transition-graph structure. Clusters are counted separately in each retained window. Degree statistics count links to clusters in adjacent windows; strength statistics sum the numbers of shared sequences on those links. Component sizes count clusters rather than member sequences.}\label{tab:transition_graph_summary}
	\begin{thesistablebody}[\footnotesize]{@{}P{0.8\textwidth}P{0.1\textwidth}@{}}
		\toprule
		\textbf{Property} & \textbf{Value} \\
		\midrule
		Clusters~\dotfill & 277,154 \\
		Directed edges~\dotfill & 117,583 \\
		Total edge weight~\dotfill & 161,116 \\
		Weak components~\dotfill & 159,693 \\
		Components containing one cluster~\dotfill & 66,879 \\
		Largest component size (clusters)~\dotfill & 860 \\
		\addlinespace
		Mean in-/out-degree~\dotfill & 0.42 \\
		Mean in-/out-strength~\dotfill & 0.58 \\
		Maximum in-degree~\dotfill & 224 \\
		Maximum out-degree~\dotfill & 178 \\
		Maximum in-strength~\dotfill & 676 \\
		Maximum out-strength~\dotfill & 496 \\
		\bottomrule
	\end{thesistablebody}
\end{table}

\begin{figure}[htbp]
	\centering
	\includegraphics[width=\textwidth]{fig_transition_graph_characteristics}
	\caption[Structure of the transition graph]{
		\textbf{Structure of the transition graph.} (A) Counts of clusters and outgoing links by retained window. (B) Complementary cumulative distribution of component size, counting the clusters connected when link directions are ignored. (C) Counts of clusters by graph role over time. (D) Total cluster counts by number of earlier and later links, using the same role colours as Panel C. Figure~\ref{fig:transition_graph_roles} illustrates these roles.
        }\label{fig:transition_graph_characteristics}
\end{figure}

The graph was fragmented, with weak components (Glossary above) dominated by isolated clusters and short chains (Figure~\ref{fig:transition_graph_characteristics}B--D). Larger temporal structures were present but uncommon. Graph-role counts showed the same pattern: simple source, sink, and isolated roles were most common, while branching and merging configurations were rare.

\subsection{Screening outcomes, null calibration, and sensitivity analysis}\label{sec:ch5_screening_outcomes}

We screened clusters using the local burst and onward burden scores, examining alternative minimum sizes and upper-tail thresholds. The primary specification, $n_v \geq 6$ and $p_{v,k}^{\mathrm{rand}} \leq 0.05$ for at least one available score, lay near the point where candidate counts began to fall more slowly as the minimum size increased (Figure~\ref{fig:threshold_robustness}B). At $p\leq0.05$, minimum sizes of four and eight gave Jaccard similarities of 0.492 and 0.607 with the primary candidate set, respectively (Figure~\ref{fig:threshold_robustness}C). These comparisons describe sensitivity of selection; they do not establish which threshold best identifies epidemiologically confirmed events.

\begin{figure}[htbp]\centering
	\includegraphics[width=\textwidth]{fig_threshold_robustness}
	\caption[How candidate selection changes with screening thresholds]{
		\textbf{How candidate selection changes with screening thresholds.} Heat maps show how the number of clusters assigned candidate status (A) and their agreement with the primary candidate-status set (C) vary across alternative minimum cluster-size thresholds and upper-tail significance thresholds. Line plots show the corresponding candidate count (B) and Jaccard similarity (D) across minimum-size thresholds at the primary upper-tail threshold, $p \leq 0.05$. The primary specification uses a minimum cluster size of 6 and $p \leq 0.05$, yielding 631 candidate-status clusters. Agreement is measured as the Jaccard similarity between the candidate-status set under each alternative specification and the primary set; red markers denote the primary specification.
        }\label{fig:threshold_robustness}
\end{figure}


Under the primary specification, the size threshold reduced the graph to 8,962 screened clusters (3.2\% of all clusters; Figure~\ref{fig:score_landscape}A). We identified 631 high-priority candidates: 428 on the local burst dimension only, 191 on the onward burden dimension only, and 12 on both dimensions. These candidates represented 7.0\% of screened clusters and 0.2\% of all transition-graph clusters.

\begin{figure}[htbp]
	\centering
	\includegraphics[width=\textwidth]{fig_score_landscape}
	\caption[Cluster eligibility, screening tiers, and score distributions]{
		\textbf{Cluster eligibility, screening tiers, and score distributions.} (A) Filtering through screening and screening-tier assignment (see Table~\ref{tab:candidate_tiers} for tier definitions). (B) Screened clusters positioned by local burst and onward burden null-standardised scores. Clusters without applicable onward burden calibration are shown in the lower strip. Point size scales with cluster size; colour indicates screening tier.
        }\label{fig:score_landscape}
\end{figure}

\begin{table}[htbp]
	\centering
	\caption[Null calibration summary for screening scores]{
    Permutation-reference summaries for screening scores. \textit{Tested $n$} counts clusters with an available calibrated score for that dimension. $p \leq 0.05$ gives the number meeting the upper-tail threshold, and \textit{Rate} gives the corresponding percentage. $p = 1$ counts values at the upper endpoint. The Kolmogorov--Smirnov (KS)\nomenclature{KS}{Kolmogorov--Smirnov} distance describes the departure of the observed $p$-value distribution from a uniform reference; a small distance does not independently establish the validity of the permutation assumptions. Conservative values include all exact ties in the upper tail; randomised values include a randomly chosen fraction of ties and were used for classification.
    }\label{tab:null_calibration_summary}
	\begin{thesistablebody}[\footnotesize]{@{}P{0.16\textwidth}P{0.18\textwidth}P{0.12\textwidth}P{0.13\textwidth}P{0.14\textwidth}P{0.12\textwidth}P{0.12\textwidth}@{}}
		\toprule
		\textbf{Dimension} & \textbf{$p$-value} & \textbf{Tested $n$} & \textbf{$p \leq 0.05$} & \textbf{Rate} & \textbf{$p = 1$} & \textbf{KS distance} \\
		\midrule
		Local burst & Conservative & 8,962 & 388 & 4.33\% & 784 & 0.087 \\
		Local burst & Randomised & 8,962 & 440 & 4.91\% & 0 & 0.004 \\
		Onward burden & Conservative & 4,216 & 195 & 4.63\% & 41 & 0.013 \\
		Onward burden & Randomised & 4,216 & 203 & 4.81\% & 2 & 0.010 \\
		\bottomrule
	\end{thesistablebody}
\end{table}

The two scores selected largely different clusters (Figure~\ref{fig:score_landscape}B). Among clusters with both null-standardised scores available, their correlation was weakly negative ($-0.20$). This describes the relationship between the constructed scores rather than demonstrating separate transmission mechanisms. Dimension-specific candidate timing and representative linked neighbourhoods are shown in Appendix~Figures~\ref{fig:app_candidate_timeline} and~\ref{fig:app_candidate_exemplars}.

Randomised upper-tail values were at most 0.05 for 4.91\% of clusters screened for local burst and 4.81\% of those screened for onward burden (Table~\ref{tab:null_calibration_summary}). Randomised tie handling substantially reduced the number of values equal to one while retaining all clusters meeting the conservative threshold. Appendix~Figure~\ref{fig:app_sse_score_null_calibration} shows the distributions. These near-nominal rates describe the observed-profile reference and tie handling; they do not estimate the fraction of genuine superspreading events or validate classification against outbreak labels.

\begin{figure}[htbp]
	\centering
	\includegraphics[width=\textwidth]{fig_stratified_calibration}
	\caption[Stratified calibration of screening-score $p$-values]{
        \textbf{Screening-score permutation $p$-values by cluster size and sequencing coverage.} Empirical cumulative distributions are shown for local burst (A--B) and onward burden (C--D). Curves are grouped by cluster size (A, C) and quartiles of window sequencing coverage calculated across screened clusters (B, D). Each cluster contributes its window's coverage, so windows with more screened clusters contribute more values to the quartile boundaries. The dashed diagonal is a uniform reference: curves above it contain a larger proportion of small values and curves below it a smaller proportion. Proximity to this line is descriptive and does not validate the screen or exclude sampling bias. Local burst panels include all screened clusters; onward burden panels include only clusters with positive observed downstream burden.
}\label{fig:stratified_calibration}
\end{figure}

The observed $p$-value distributions were similar across coverage quartiles and close to the uniform reference (Figure~\ref{fig:stratified_calibration}). Differences by cluster size were clearer: small local burst $p$-values were more common among larger clusters, while onward burden showed a weaker pattern in the opposite direction. These associations are consistent with including relative size in local burst and dividing burden by starting size before ranking. Neither similarity across coverage groups nor these size patterns establishes robustness to unobserved infections or targeted sampling.

\section{Discussion}\label{sec:ch5_discussion}

The practical problem was to prioritise clusters for investigation when large sampled groups can arise under different surveillance conditions. This chapter provided an explicit comparison rule using relative size, newly represented membership, and additional sequences in later linked clusters. Of 277,154 clusters across retained windows, 8,962 met the minimum size and 631 received CSC status: 7.0\% of screened clusters and 0.23\% of all clusters. The screen supplies a defined review set whose underlying memberships and score components can be examined.

\subsection{Cluster patterns reflect transmission and observation}\label{sec:ch5_cluster_patterns}

Most clusters were small and a few were very large. Studies of individual transmission and viral introductions provide reasons to expect substantial variation in outbreak size \citep{LloydSmith2005,Endo2020,daSilvaFilipe2020,duPlessis2021,GomezCarballa2021}. Here, however, the observed distribution also depends on sequencing, window boundaries, and the partition of compatibility graphs. Its skew alone cannot identify the contribution of transmission heterogeneity. Similarly, the median within-cluster distances of 0 SNPs and 4 days describe a partition constructed using genetic and temporal information; they are not an independent test of its transmission interpretation.

Large clusters were most common during periods of high recorded incidence and sequencing volume, especially Omicron waves. Chapter~\ref{ch:genomic_networks} showed that sequence volume, case counts, and sequencing coverage had different temporal patterns. Their associations with cluster summaries motivate comparison within a sampling context, but do not separate effects of transmission from ascertainment. Existing methods that infer reproduction numbers or transmission heterogeneity from genomic cluster sizes require an explicit account of how transmission and sampling produce their clusters \citep{TranKiem2024}. Methods following persistent genomic groups also differ from the repeated community assignments used here \citep{dor2025}. The present screen does not estimate those epidemic parameters.

Typical clusters contained records from few residential Data Zones. Larger clusters represented more zones without necessarily having greater median residential distances. This distinguishes the number of areas represented from the distance between them. Chapter~\ref{ch:genomic_networks} provides complementary evidence of compatibility within administrative areas and urban/rural classes, but neither shared administrative membership nor a shared urban/rural class specifies a transmission setting. Residential concentration can guide follow-up questions about local exposures without identifying those exposures.

\subsection{Calibrated screening distinguishes local bursts from onward burden}\label{sec:ch5_discussion_screening}

The screen asks whether a cluster has a high score relative to its observed comparison group. Ranking within windows and choosing comparison groups by available context make that reference explicit. They do not reconstruct missing infections or demonstrate that changes in testing, sequencing coverage, or outbreak-focused sampling have been removed. The minimum size and upper-tail cutoffs determine the workload and composition of the resulting review set; agreement between nearby choices describes operational sensitivity.

Only 12 candidates met both criteria, compared with 428 meeting the local burst criterion alone and 191 meeting the onward burden criterion alone. The two dimensions therefore selected different sampled patterns. This is relevant to the distinction between a large accumulation and later wider representation motivated by the Boston investigations \citep{Lemieux2021}. However, the scores also differ by construction: local burst includes relative size, onward burden adjusts for starting size, and onward burden is available only for clusters with additional later membership. The limited overlap does not independently demonstrate two biological types of outbreak.

The transition graph makes these membership comparisons possible. Its fragmented structure describes discontinuities in observed links. Such gaps may reflect the end of a sampled group, missed infections, changes in cluster assignments, separate lineage analyses, or window boundaries. The graph cannot distinguish these explanations or establish that transmission chains were short or contained. Cumulative burden also depends on how much later observation is available: clusters near the end of the record have less opportunity for later linked sequences to be observed.

\subsection{Candidate status defines hypotheses, not events}\label{sec:ch5_strengths_limitations}

The main strength is the explicit connection between each candidate label, its observed membership, and the comparison rule. This permits epidemiological review of why a cluster was selected and supports the subsequent characterisation of its sampled population. Age, sex, deprivation, and residential attributes did not determine candidate status.

Calibration requires further methodological assessment. Reassigning complete observed profiles preserves their component relationships. When the permutation groups lie within a window, it redistributes existing composite scores among clusters; near-uniform upper-tail values can therefore arise from the ranking reference itself. Their distribution does not independently validate a null model for transmission. Inferential interpretation requires exchangeability: profiles must be legitimately interchangeable within the comparison groups under the stated null \citep{Hemerik2017}. Shared sequences, linked clusters, and overlapping windows create dependence that this grouping does not by itself resolve. For onward burden, permutation across windows within clades also requires assessment in relation to changing sampling and unequal follow-up. These questions remain unresolved by the descriptive diagnostics presented here.

Further evaluation should examine sensitivity to window design, clustering choices, component weighting, and the observation horizon, and compare selected clusters with independently documented exposures or outbreaks. Until that evaluation, CSC status is a hypothesis for investigation. Neither a small permutation value nor selection on both dimensions confirms a superspreading event.

Chapter~\ref{ch:candidate_characterisation} uses these graph-defined outcomes to ask which sampled populations, places, policy periods, and clades are represented among CSC clusters. This follows the thesis's analytical sequence: estimate compatibility between sampled infections, screen the resulting clusters for unusual patterns, and then examine their epidemiological context.

\section{Conclusion}\label{sec:ch5_discussion_conclusions}

Most clusters were small and represented few residential Data Zones, while larger clusters were more common during periods with high sequence volumes. Of 8,962 clusters in retained windows containing at least six sequences, the calibrated screen assigned CSC status to 631: 428 for local burst only, 191 for onward burden only, and 12 for both. These results define a review set based on different patterns of sampled membership; their epidemiological specificity remains to be assessed.

Chapter~\ref{ch:candidate_characterisation} examines how population composition and within-cluster diversity relate to CSC status and the two screening scores.
